\documentclass[9pt]{beamer}
\usetheme{Madrid}

\usepackage{polyglossia}
\setdefaultlanguage{spanish}

\usepackage{fontspec}
% switch to a monospace font supporting more Unicode characters
\setmonofont{JetBrains Mono}[
    Contextuals=Alternate,
    Ligatures=TeX,
    Scale=MatchLowercase,
]

\usepackage{minted}
\usemintedstyle{pastie}

% === Math facilities ===

% Theorem environments
\usepackage{amsmath}
\usepackage{amsthm}
\usepackage{amssymb}
\usepackage{centernot}
\usepackage{mathtools}

\newcommand{\funfont}[1]{\texttt{#1}}
\newcommand{\tyfont}[1]{\texttt{\textbf{#1}}}
\renewcommand{\implies}{\rightarrow}
\renewcommand{\impliedby}{\leftarrow}
\renewcommand{\iff}{\leftrightarrow}
\renewcommand{\mapsto}{\Rightarrow}
\newcommand{\Prop}{\mathbf{Prop}}
\newcommand{\placeholder}{\square}

\newcommand{\ImpIfElse}[3]{\texttt{if} \ {#1} \ \{{#2}\} \ \texttt{else} \ \{{#3}\}}
\newcommand{\ImpWhile}[2]{\texttt{while} \ {#1} \ \{{#2}\}}
\newcommand{\ImpConcat}[2]{{#1} \texttt{;} \ {#2}}
\newcommand{\ImpSkip}{\texttt{skip}}
\newcommand{\ImpAssign}[2]{{#1} = {#2}}

\newcommand{\LeanIfElse}[3]{\text{if} \ {#1} \ \text{then} \ {#2} \ \text{else} \ {#3}}
\newcommand{\LeanTrue}{\text{true}}
\newcommand{\LeanFalse}{\text{false}}

\newcommand{\ImpPlus}{\mathbin{\texttt{+}}}
\newcommand{\ImpMinus}{\mathbin{\texttt{-}}}
\newcommand{\ImpTimes}{\mathbin{\texttt{*}}}

\newcommand{\ImpNot}{\texttt{!}}
\newcommand{\ImpAnd}{\mathbin{\texttt{\&\&}}}
\newcommand{\ImpOr}{\mathbin{\texttt{||}}}
\newcommand{\ImpEq}{\mathbin{\texttt{==}}}
\newcommand{\ImpLe}{\mathbin{\texttt{<=}}}

\newcommand{\ImpTrue}{\texttt{true}}
\newcommand{\ImpFalse}{\texttt{false}}

\newcommand{\ImpName}{\texttt{IMP}}

\usepackage[only,llbracket,rrbracket]{stmaryrd}

% define math symbol for big step arrow
\newcommand{\bigstep}{\Longrightarrow}
\newcommand{\Denote}[1]{\llbracket{} {#1} \rrbracket{}}

% === End math facilities ===

\usepackage{csquotes}

%Information to be included in the title page:
\title{Presentación con Lean 4 de algunos Resultados clásicos sobre
      Semántica de Lenguajes de Programación}
\author{Ricardo Maurizio Paul}
% \date{2021}

\begin{document}

\frame{\titlepage}

\begin{frame}
\frametitle{Índice}
\tableofcontents
\end{frame}

\section{Introducción}
\begin{frame}
\frametitle{Motivación}
\begin{itemize}
  \item Los asistentes de demostración combinan las matemáticas clásicas con la programación.
  \item La semántica de lenguajes de programación es un campo particularmente adecuado
    para explorar las complicaciones que entraña el uso de asistentes de demostración.
  \item Lean 4 es, además, un lenguaje de programación funcional. Tener demostraciones formales
    sobre la semántica de un lenguaje nos facilitaría, por ejemplo, crear
    un compilador verificado para dicho lenguaje. En particular estas semánticas son
    especialmente útiles para demostrar la corrección de las diferentes transformaciones
    que se pueden aplicar a un programa.
\end{itemize}

\end{frame}

\begin{frame}
\frametitle{Objetivos}
\begin{itemize}
  \item Formalizar en Lean 4 algunos de los resultados vistos en la asignatura de Teoría de la Programación,
    en particular definir la semántica natural, estructural y denotacional.
  \item Demostrar la equivalencia entre las semánticas, lo que nos permite
    utilizar la semántica que más nos convenga en cada caso.
  \item Intentar que el código sea lo más legible posible, evitando generalizaciones
    innecesarias y el \emph{simplificador}, todas las demostraciones son \enquote{elementales}. 
  \item Que el código sea completo, por lo que hemos evitado \emph{Mathlib}.
\end{itemize}

\end{frame}

\begin{frame}
\frametitle{el lenguaje \ImpName{}}
Definimos la sintaxis abstracta de \(\tyfont{Aexp}\), en forma de Backus-Naur:
\[
a \Coloneqq n \ | \ x \ | \ a_1 \ImpPlus a_2 \ | \ a_1 \ImpMinus a_2 \ | \ a_1 \ImpTimes a_2
\]
la sintaxis de \(\tyfont{Bexp}\) como:
\[
b \Coloneqq \ImpTrue \ | \ \ImpFalse \ | \ 
  \ImpNot b \ | \ b_1 \ \ImpAnd \ b_2 \ | \ b_1 \ \ImpOr \ b_2 \ | \ a_1 \ \ImpEq \ a_2 \ | \ a_1 \ \ImpLe \ a_2
\]
y la sintaxis de \(\tyfont{Com}\) como:
\[
c \Coloneqq \ImpSkip \ | \ \ImpAssign{x}{a} \ | \ \ImpConcat{c_1}{c_2} \ | \ \ImpIfElse{b}{c_1}{c_2} \ | \ \ImpWhile{b}{c}
\]
\end{frame}

\section{Semántica Operacional}
\begin{frame}
\frametitle{Semántica natural}

La definición de la semántica de paso largo es casi igual que en~\cite{winskel_formal_semantics}.
Las \enquote{reglas} de la traducción dirigida por sintaxis se formalizan mediante constructores.
De esta forma podemos \enquote{reutilizar} gran parte de las demostraciones, mediante
la inducción estructural sobre los constructores.

\end{frame}

\begin{frame}
\frametitle{Constructores}
La semántica de paso largo es una relación \(\tyfont{Com} \times \tyfont{State} \implies \tyfont{State} \implies \Prop\)
con los siguientes constructores (\enquote{reglas}):
\begin{align*}
  \ImpSkip &: (\ImpSkip, s) \bigstep s \\
  \funfont{ass} &: (v = a, s) \bigstep s[v \leftarrow a \ s] \\
  \funfont{cat} &: \frac{
    (c_1, s_1) \bigstep s_2 \quad (c_2, s_2) \bigstep s_3
  }{
    (\ImpConcat{c_1}{c_2}, s_1) \bigstep s_3
  } \\
  \funfont{ifFalse} &: \frac{b \ s_1 = \LeanFalse \quad (c_2, s_1) \bigstep s_2}
    {(\ImpIfElse{b}{c_1}{c_2}, s_1) \bigstep s_2} \\
  \funfont{ifTrue} &: \frac{b \ s_1 = \LeanTrue \quad (c_1, s_1) \bigstep s_2}
{(\ImpIfElse{b}{c_1}{c_2}, s_1) \bigstep s_2} \\
\funfont{whileFalse} &: \frac{b \ s = \LeanFalse}
  {(\ImpWhile{b}{c}, s) \bigstep s} \\
  \funfont{whileTrue} &: \frac{b \ s_1 = \LeanTrue \quad (c, s_1) \bigstep s_2 \quad
  (\ImpWhile{b}{c}, s_2) \bigstep s_3}
{(\ImpWhile{b}{c}, s_1) \bigstep s_3}
\end{align*}
\end{frame}

\begin{frame}[fragile]
\frametitle{Definición en Lean}
\begin{minted}[escapeinside={}]{lean4}
inductive BigStep: Com × State → State → Prop
  | skip: BigStep (∅, s) s
  | ass: BigStep (ass v a, s) (s[v ← a s])
  | cat (s'': State) (hcatl: BigStep (c, s) s'') 
        (hcatr: BigStep (c', s'') s'):
    BigStep (c++c', s) s'
  | ifFalse (hcond: b s = false) 
        (hifright: BigStep (c', s) s'):
    BigStep (ifElse b _ c', s) s'
  | ifTrue (hcond: b s) (hifleft: BigStep (c, s) s'):
    BigStep (ifElse b c _, s) s'
  | whileFalse (hcond: b s = false):
    BigStep (whileLoop b c, s) s
  | whileTrue (s'': State) (hcond: b s) 
        (hwhilestep: BigStep (c, s) s'')
        (hwhilerest: BigStep (whileLoop b c, s'') s'):
    BigStep (whileLoop b c, s) s'
\end{minted}
\end{frame}

\begin{frame}[fragile]
\frametitle{Árboles de derivación 1}
    Los \enquote{árboles de derivación} de~\cite{winskel_formal_semantics} se transforman
    entonces en la aplicación sucesiva de constructores. Consideramos el siguiente ejemplo:
\begin{minted}[escapeinside={}]{rust}
x = 2;
if x <= 1 {
  y = 3
} else {
  z = 4
}
\end{minted}
\end{frame}

\begin{frame}
\frametitle{Árboles de derivación 2}
Para aligerar la notación denotaremos el programa como \(P\), y la estructura 
\texttt{if} como \(\text{Ifelse}\), entonces:
\[
F \coloneqq \funfont{ifFalse}: \displaystyle{\frac{
    \funfont{ass}: (\ImpAssign{z}{4} , s_0[x \leftarrow 2]) \bigstep s_0[x \leftarrow 2][z \leftarrow 4]
    }
{
    (\text{Ifelse}, s_0[x \leftarrow 2]) \bigstep s_0[x \leftarrow 2][z \leftarrow 4]
}}
\]
y el árbol completo es:
\[
\funfont{cat}: \frac{
\funfont{ass}: (\ImpAssign{x}{2}, s_0) \bigstep s_0[x \leftarrow 2]
\qquad
F
}{
(P, s_0) \bigstep s_0[x \leftarrow 2][z \leftarrow 4]
}
\]
\end{frame}

\begin{frame}[fragile]
\frametitle{Ejemplo de árbol de derivación en Lean}
\begin{minted}[escapeinside={}]{lean4}
private example: 
  ([|x = 2; y = 1; while 2 <= x {y = y * x; x = x - 1}|], s0)
    ==> s0["x" ← 2]["y" ← 1]["y" ← 2]["x" ← 1] 
  := by
  apply cat (s0["x" ← 2]["y" ← 1]) (cat (s0["x" ← 2]) ass ass)
  apply whileTrue (s0["x" ← 2]["y" ← 1]["y" ← 2]["x" ← 1])
  apply rfl
  apply cat (s0["x" ← 2]["y" ← 1]["y" ← 2]) ass ass
  apply whileFalse
  apply rfl
\end{minted}
\end{frame}

\begin{frame}[fragile]
\frametitle{Reglas de reescritura}

Vemos como podemos realizar una demostración usando la semántica de paso largo.

\begin{minted}[escapeinside={}]{lean4}
theorem if_eq': 
  (ifElse b c c', s) ==> s' ↔ (cond (b s) c c', s) ==> s' 
  := {
    mp := fun hmp =>  match hmp with 
      | ifTrue hc h 
      | ifFalse hc h => hc ▸ h,
    mpr := match hb: b s with 
      | true => (ifTrue hb .) 
      | false => (ifFalse hb .)
  }
\end{minted}
\end{frame}

\begin{frame}[fragile]
\frametitle{Reglas de reescritura 2}

\begin{minted}[escapeinside={}]{lean4}
theorem while_eq: (whileLoop b c, s) ==> s' ↔
  cond (b s) (∃s'', (c, s) ==> s'' ∧ (whileLoop b c, s'') ==> s') 
    (s = s') := {
    mp := fun hmp => match hmp with
      | whileTrue s'' hc h hw => hc ▸ ⟨s'', h, hw⟩
      | whileFalse hb => hb ▸ rfl,
    mpr := match hb: b s with
      | true => fun ⟨s'', h, hw⟩ => whileTrue s'' hb h hw
      | false => (. ▸ whileFalse hb)
  }
\end{minted}

\end{frame}

\begin{frame}[fragile]
\frametitle{Desenrollado de bucles}

\begin{minted}[escapeinside={}]{lean4}
instance equiv: Setoid Com where
  r c c' := ∀{s s': State}, (c, s) ==> s' ↔ (c', s) ==> s'
  iseqv := ⟨fun _ => Iff.rfl, (Iff.symm .), (Iff.trans . .)⟩

theorem loop_unfold: whileLoop b c ≈ ifElse b (c++whileLoop b c) ∅ 
  := by
  intro s _; rw [if_eq', while_eq]
  match b s with
  | false => exact skip_eq.symm
  | true => exact ⟨fun ⟨_, hl, hr⟩ => cat _ hl hr,
    fun (cat s'' hl hr) => ⟨s'', hl, hr⟩⟩
\end{minted}

\end{frame}

\section{Semántica estructural}

\begin{frame}
\frametitle{Semántica estructural}

Podemos definir la semántica estructural de forma similar.
\[
\begin{gathered}
  \funfont{ass}: (\ImpAssign{v}{a}, s) \rightsquigarrow (\ImpSkip, s[v \leftarrow s]) \\
  \funfont{skipCat}: (\ImpConcat{\ImpSkip}{c}, s) \rightsquigarrow (c, s) \qquad
  \funfont{cat}: \frac{(c, s) \rightsquigarrow (c', s')}{(\ImpConcat{c}{c''}, s) \rightsquigarrow (c', c'', s')} \\
  \funfont{ifElse}: (\ImpIfElse{b}{c}{c'}, s)
    \rightsquigarrow (\LeanIfElse{b \ s}{c}{c'}, s) \\
  \funfont{whileLoop}: (\ImpWhile{b}{c}, s)
    \rightsquigarrow (\LeanIfElse{b \ s}{\ImpConcat{c}{\ImpWhile{b}{c}}}{\ImpSkip}, s) 
\end{gathered}
\]

Para poder dar \enquote{más de un paso} tendremos que definir también la 
\emph{clausura reflexiva transitiva} de la relación. Las demostraciones
sobre esta clausura se harán mediante inducción en la \enquote{cabeza} o
\enquote{cola} de la relación.

\end{frame}

\begin{frame}[fragile]
\frametitle{Definición en Lean}
\begin{minted}[escapeinside={}]{lean4}
inductive SmallStep: Com × State → Com × State → Prop where
  | ass: SmallStep (ass v a, s) (∅, s[v ← a s])
  | skipCat: SmallStep (∅++c, s) (c, s)
  | cat (hstep: SmallStep (c, s) (c', s')):
    SmallStep (c++c'', s) (c'++c'', s')
  | ifElse: SmallStep (ifElse b c c', s) (cond (b s) c c', s)
  | whileLoop:
    SmallStep (whileLoop b c, s) (cond (b s) (c++whileLoop b c) ∅, s)
\end{minted}

\end{frame}

\begin{frame}
\frametitle{Clausura reflexiva transitiva}
La clausura reflexiva transitiva de una relación binaria \(\mathcal{R}\) es la
  relación más pequeña que contiene \(\mathcal{R}\) y es tanto reflexiva como
  transitiva. En concreto es un tipo inductivo 
  \(\tyfont{ReflTrans}: (\alpha \to \alpha \to \Prop) \to
  (\alpha \to \alpha \to \Prop)\) y la denotamos como
  \(\mathcal{R}^* \coloneqq \tyfont{ReflTrans} \ \mathcal{R}\).

  El tipo \(\tyfont{ReflTrans}\) tiene dos constructores:
  \[
    \funfont{refl}: a \mathcal{R}^* a
    \qquad 
    \funfont{tail}: \frac{
      a \mathcal{R}^* b
      \quad b \mathcal{R} c
    }{a \mathcal{R}^* c}
  \]
\end{frame}

\begin{frame}
\frametitle{Árbol de derivación}
Consideramos el programa \texttt{x = 1; while x <= 1 \{x = x + 1\}},
  para facilitar la lectura denotaremos el bloque \enquote{while} como \(W\).
  Un árbol de derivación parcial es:
  \[
    H \coloneqq
    \funfont{head}: \displaystyle{\frac{
        \funfont{skipcat}: (\ImpConcat{\ImpSkip}{W}, s_0[x \leftarrow 1]) \rightsquigarrow^* (W, s_0[x \leftarrow 1])
        \qquad
        \cdots
      }{
        (\ImpConcat{\ImpSkip}{W}, s_0[x \leftarrow 1]) \rightsquigarrow^* (W, s_0[x \leftarrow 1][x \leftarrow 2])
      }}
  \]
  y por tanto:    
  \[
    \funfont{head}: \frac{
      \funfont{cat}: \displaystyle{\frac{
        \funfont{ass}: (\ImpAssign{x}{1}, s_0) \rightsquigarrow (\ImpSkip, s_0[x \leftarrow 1])
      }
      {(\ImpConcat{\ImpAssign{x}{1}}{W}, s_0)
      \rightsquigarrow
      (\ImpConcat{\ImpSkip}{W}, s_0[x \leftarrow 1])
      }}
      \qquad
      H
    }{
      (\ImpConcat{\ImpAssign{x}{1}}{W}, s_0) \rightsquigarrow^* (W, s_0[x \leftarrow 1][x \leftarrow 2])
    }
  \]
  que hemos obtenido aplicando sucesivamente los constructores. Notamos que
  dando \enquote{un paso más} llegaríamos al estado final del programa.
\end{frame}

\section{Semántica denotacional relacional}

\begin{frame}
\frametitle{Semántica denotacional relacional}

La semántica denotacional necesita \emph{funciones parciales}, ya que hay programas que no acaban.

Para poder desarrollar la semántica denotacional de \ImpName{} como en~\cite{winskel_formal_semantics},
introducimos los conjuntos y la teoría de orden necesaria sobre los mismos. Además demostramos
los teoremas de punto fijo de Knaster-Tarski y Kleene. 
Añadimos a esta semántica el
adjetivo \enquote{relacional} ya que analizaremos \emph{relaciones} en vez de \emph{funciones},
este enfoque es tomado de~\cite{semantics_with_isabelle}.
\end{frame}

\begin{frame}[fragile]
\frametitle{Semántica denotacional}
\begin{minted}[escapeinside={}]{lean4}
def eval (c: Com) (s: State): Option State :=
  match c with
  | skip => some s
  | ass v a => some (s[v ← a s])
  | cat c c' => eval c s >>= eval c'
  | ifElse b c c' => eval (if b s then c else c') s
  | whileLoop b c' => if b s then eval c' s >>= eval c else s
partial_fixpoint
\end{minted}
\end{frame}


\begin{frame}[fragile]
\frametitle{Funciones como relaciones}
\begin{minted}[escapeinside={}]{lean4}
def id: Set (α × α) := {p | p.1 = p.2}

theorem mem_id: (a, b) ∈ id ↔ a = b :=
  Iff.rfl

def comp (f g: Set (α × α)): Set (α × α) :=
  {(x, z) | ∃y, (x, y) ∈ f ∧ (y, z) ∈ g}

infixl:90 " ○ " => SRel.comp

theorem mem_comp: x ∈ f ○ g ↔ ∃z, (x.1, z) ∈ f ∧ (z, x.2) ∈ g :=
  Iff.rfl
\end{minted}
\end{frame}

\begin{frame}[fragile]
\frametitle{Semántica denotacional}
\begin{minted}[escapeinside={}]{lean4}
private def W (b: Bexp) (f: Set (State × State)):
  Set (State × State) → Set (State × State) :=
  fun g => {(s, _) | b s}.ite (f ○ g) SRel.id
  
instance W.OrderHom (b: Bexp) (f: Set (State × State)):
  Set (State × State) →o Set (State × State) :=
    ⟨W b f, W.Monotone⟩

def denote: Com → Set (State × State)
  | skip => SRel.id
  | ass v a => {(s, t) | t = s[v ← a s]}
  | cat c c' => c.denote ○ c'.denote
  | ifElse b c c' => {(s, _) | b s}.ite c.denote c'.denote
  | whileLoop b c => (W.OrderHom b c.denote).lfp
\end{minted}
\end{frame}

\section{Equivalencia}


\begin{frame}[fragile]
\frametitle{Equivalencia}

La demostración de equivalencia se realiza en dos pasos.
\begin{itemize}
    \item Demostramos que la semántica natural implica la estructural y viceversa.
    \item Demostramos que la semántica natural implica la denotacional y viceversa.
\end{itemize}

\end{frame}

\begin{frame}[fragile]
\frametitle{Estructural implica natural}
\begin{minted}[escapeinside={}]{lean4}
theorem BigStep.from_structural_step
  (h: x ~> x') (h': x' ==> s'): x ==> s' := by
  induction h generalizing s' with
  | ass => exact (skip_eq.mp h') ▸ ass
  | skipCat => exact cat _ skip h'
  | cat _ ih => match h' with
    | cat s' hcatl hccatr => exact cat s' (ih hcatl) hccatr
  | ifElse => rw [if_eq']; exact h'
  | whileLoop => rw [loop_unfold, if_eq']; exact h'

theorem BigStep.from_structural (h: x ~>* (∅, s')): x ==> s' := by
  induction h using ReflTrans.head_induction_on with
  | refl => exact skip
  | head hsingle _ hs' => exact from_structural_step hsingle hs'
\end{minted}
\end{frame}

\begin{frame}[fragile]
\frametitle{Denotacional implica natural}
\begin{minted}[escapeinside={}]{lean4}
theorem OrderHom.lfp_le [CompleteLattice α] {f: α →o α} (h: pfp f a):
  lfp f ≤ a := CompleteLattice.Inf_le _ h

theorem BigStep.from_denote (h: (s, s'') ∈ ⟦c⟧): (c, s) ==> s'' := by
  induction c generalizing s s'' with
  | skip => exact SRel.mem_id.mp h ▸ skip
  | ass => exact h ▸ ass
  | cat _ _ ihcatl ihcatr =>
    exact match h with | ⟨s', hl, hr⟩ => cat s' (ihcatl hl) (ihcatr hr)
  | ifElse _ _ _ ih1 ih2 =>
    exact match h with
    | .inl ⟨hs, hc⟩ => ifTrue hc (ih1 hs)
    | .inr ⟨hs, hc⟩ => ifFalse (Bool.not_eq_true _ ▸ hc) (ih2 hs)
  | whileLoop b c ih =>
    have hle: ⟦whileLoop b c⟧ ≤ {(s, s'') | (whileLoop b c, s) ==> s''} :=
      OrderHom.lfp_le fun (_, _) hmp =>
        match hmp with
        | .inl ⟨⟨s', hs, hrest⟩, hc⟩ =>
          whileTrue s' hc (ih hs) hrest
        | .inr ⟨hid, hc⟩ =>
          (SRel.mem_id.mp hid) ▸ whileFalse (Bool.not_eq_true _ ▸ hc)

    exact hle h
\end{minted}
\end{frame}

\section{Bibliografía}

\begin{frame}
\frametitle{Bibliografía}
\bibliographystyle{alpha}
\bibliography{refs}
\end{frame}

\end{document}